\section{集合的含义与表示}

\begin{frame}{集合的概念}
  一般地，我们把研究对象统称为\emph{元素}，把一些元素组成的总体叫做\emph{集合}。

  \begin{block}{集合中元素的三个特性}
    \begin{enumerate}
      \item 确定性：元素是否属于集合是明确的；
      \item 互异性：集合中的元素互不相同；
      \item 无序性：元素之间没有顺序之分。
    \end{enumerate}
  \end{block}
\end{frame}

\begin{frame}{元素与集合的关系}
  若 $a$ 是集合 $A$ 的元素，称 $a$ \alert{属于} $A$，记作 $a\in A$；
  否则称 $a$ \alert{不属于} $A$，记作 $a\notin A$。

  \vspace{1em}
  \begin{exampleblock}{常用数集}
    \[
      \mathbb{N}\ (\text{自然数集}),\quad
      \mathbb{Z}\ (\text{整数集}),\quad
      \mathbb{Q}\ (\text{有理数集}),\quad
      \mathbb{R}\ (\text{实数集}).
    \]
  \end{exampleblock}
\end{frame}

\begin{frame}{集合的表示方法}
  \begin{columns}[T]
    \begin{column}{0.48\textwidth}
      \textbf{列举法}
      \[
        \{1,2,3,4,5\}
      \]
      把集合中的元素一一列举出来。
    \end{column}
    \begin{column}{0.48\textwidth}
      \textbf{描述法}
      \[
        \{x\in\mathbb{R}\mid x^2-1=0\}
      \]
      用元素的共同特征描述集合。
    \end{column}
  \end{columns}
\end{frame}

\begin{frame}{例题}
  \begin{block}{例}
    用列举法表示集合 $\{x\in\mathbb{Z}\mid -2<x\le 3\}$。
  \end{block}
  \pause
  逐一检验整数：$-1,0,1,2,3$，故
  \[
    \{x\in\mathbb{Z}\mid -2<x\le 3\}=\{-1,0,1,2,3\}.
  \]
\end{frame}

\begin{frame}{课堂练习}
  用适当的方法表示下列集合：
  \begin{enumerate}
    \item 由 $1,2,3,4,5$ 组成的集合；
    \item 函数 $y=x^2-1$ 的图象与 $x$ 轴交点的横坐标；
    \item 所有大于 $3$ 的整数。
  \end{enumerate}
\end{frame}
